\documentclass[11pt]{article}

\usepackage[margin=1in]{geometry}
\usepackage{amsmath,amssymb}
\usepackage{booktabs}

\title{CatBoost + CNN Residual Learning: Formal Description}
\date{}

\begin{document}
\maketitle

\section{Notation and Dataset}

Let $i \in \{1,\dots,N_c\}$ index companies and let $t$ index months. Each row in the
panel is a company--month pair $(i,t)$. The target is the gross return
$y_{i,t}$ stored in \texttt{monthly\_gross\_return}. For interpretation, define the
net return $r_{i,t} = y_{i,t} - 1$, but the pipeline models $y_{i,t}$ directly.
Let $\mathbf{x}_{i,t} \in \mathbb{R}^d$ be the raw feature vector for $(i,t)$, and
let $X \in \mathbb{R}^{N \times d}$ collect all rows, with row index
$r = r(i,t)$ so that $X_{r,:} = \mathbf{x}_{i,t}$.

The dataset is a monthly panel containing:
\begin{itemize}
	\item \textbf{Identifiers:} \texttt{co\_code}, \texttt{Month}.
	\item \textbf{Target:} \texttt{monthly\_gross\_return}.
	\item \textbf{Market variables:} \texttt{mktcap}, \texttt{price},
	\texttt{lag\_mv}, \texttt{BM\_sep}.
	\item \textbf{Accounting variables:} \texttt{equity\_bv\_on\_stkdate},
	\texttt{fy\_book\_value}, \texttt{sa\_net\_worth}, \texttt{sa\_total\_assets},
	\texttt{sa\_sales}, \texttt{sa\_cost\_of\_goods\_sold},
	\texttt{sa\_total\_interest\_exp}, \texttt{sa\_selling\_distribution\_exp},
	\texttt{lagged\_book\_equity}, \texttt{lagged\_assets}, \texttt{lagged2\_assets}.
	\item \textbf{Factor signals:} \texttt{OpProf}, \texttt{Inv}, \texttt{Momentum}.
	\item \textbf{Categorical labels:} \texttt{Size\_Label} (S/B),
	\texttt{BM\_Label} (G/N/V), \texttt{OpProf\_Label} (W/N/R),
	\texttt{Inv\_Label} (A/N/C), \texttt{Mom\_Label} (L/N/W).
\end{itemize}

Company \texttt{co\_code = 194} is removed from the initial dataset and treated as
a standalone holdout company for out-of-sample evaluation.

\section{Problem Formulation}

The learning task is to predict the gross return $y_{i,t}$ using information
available up to month $t$. We write the prediction as
\[
\hat{y}_{i,t} = f_{\theta}(\mathcal{I}_{i,t}),
\]
where $\mathcal{I}_{i,t}$ includes the current engineered features and short
historical windows. The objective is a regression loss such as mean squared
error (MSE):
\[
\theta^* = \arg\min_{\theta} \frac{1}{N} \sum_{(i,t)} (y_{i,t} - \hat{y}_{i,t})^2.
\]

\section{Preprocessing Summary}

The preprocessing stage (see preprocess.tex) ensures time ordering and
leakage-safe features by:
\begin{itemize}
	\item coercing types (dates, numeric, categorical),
	\item engineering lags and rolling statistics from past returns,
	\item adding missingness indicators,
	\item splitting into train/validation/test and masking test targets,
	\item removing raw return columns \texttt{lag\_ret} and \texttt{log\_ret}.
\end{itemize}

We denote the final tabular feature vector after preprocessing by
$\tilde{\mathbf{x}}_{i,t} \in \mathbb{R}^{d'}$.

\section{Model Structure}

The pipeline uses a two-stage additive model:
\[
\hat{y}_{i,t} = \underbrace{f_{\text{CB}}(\tilde{\mathbf{x}}_{i,t})}_{\text{CatBoost}}
\; + \; \underbrace{f_{\text{CNN}}(\mathcal{S}_{i,t})}_{\text{residual correction}}.
\]

The diagram below summarizes the flow:
\[
\begin{array}{c}
\text{processed panel } (i,t)
\\ \downarrow \\
f_{\text{CB}} \;\Rightarrow\; \hat{y}^{\text{CB}}_{i,t}
\\ \downarrow \\
e_{i,t} = y_{i,t} - \hat{y}^{\text{CB}}_{i,t}
\\ \downarrow \\
f_{\text{CNN}} \;\Rightarrow\; \hat{e}_{i,t}
\\ \downarrow \\
\hat{y}_{i,t} = \hat{y}^{\text{CB}}_{i,t} + \hat{e}_{i,t}
\end{array}
\]

\section{CatBoost Stage}

CatBoost is trained on the train split with categorical columns passed
explicitly. The model minimizes RMSE on the training set and uses the validation
split for early stopping.

\subsection{CatBoost Residuals}

After fitting CatBoost, the residual for each row is defined as
\[
e_{i,t} = y_{i,t} - \hat{y}^{\text{CB}}_{i,t},
\quad \hat{y}^{\text{CB}}_{i,t} = f_{\text{CB}}(\tilde{\mathbf{x}}_{i,t}).
\]
Residuals on the validation split become the training targets for the CNN.

\section{Residual CNN Stage}

\subsection{CNN Input Construction}

Let $L=6$ be the sequence length. For each company $i$ and month $t$ with at
least $L$ months of history, define the sequence
\[
\mathcal{S}_{i,t} =
\begin{bmatrix}
\mathbf{z}_{i,t-L+1}^\top \\
\vdots \\
\mathbf{z}_{i,t}^\top
\end{bmatrix}
\in \mathbb{R}^{L \times d'},
\]
where $\mathbf{z}_{i,t}$ is the normalized feature vector formed by concatenating
$\tilde{\mathbf{x}}_{i,t}$ with the CatBoost prediction
$\hat{y}^{\text{CB}}_{i,t}$. Each numeric component is standardized using
\[
\mathbf{z}_{i,t} = \frac{\tilde{\mathbf{x}}_{i,t} - \boldsymbol{\mu}}{\boldsymbol{\sigma}},
\]
with $\boldsymbol{\mu}$ and $\boldsymbol{\sigma}$ computed on the validation
sub-split used for CNN training. Non-finite values are set to zero after
normalization. Categorical labels are encoded as integer category codes using a
global category map; unseen categories map to \texttt{\_\_MISSING\_\_}.

\subsection{CNN Architecture}

The residual CNN is a 1D convolutional network over time:
\begin{itemize}
	\item input channels: $d'$ (features plus \texttt{cb\_pred}),
	\item Conv1d(64, kernel=3, padding=1) + ReLU,
	\item Conv1d(32, kernel=3, padding=1) + ReLU,
	\item AdaptiveAvgPool1d(1),
	\item Linear(32 $\to$ 1).
\end{itemize}
The network outputs $\hat{e}_{i,t}$, the predicted residual.

\subsection{CNN Objective}

The CNN minimizes MSE on residuals:
\[
\min_{\phi} \frac{1}{N_{\text{seq}}} \sum_{(i,t)}
\left(e_{i,t} - f_{\text{CNN}}(\mathcal{S}_{i,t})\right)^2.
\]

\section{Training Protocol}

\begin{itemize}
	\item \textbf{CatBoost:} trained on \texttt{train} with validation for early
	stopping; GPU is used when available, with CPU fallback.
	\item \textbf{CNN:} trained on the validation split only
	(\texttt{USE\_FULL\_VALIDATION\_FOR\_CNN = True}).
	\item \textbf{Random seeds:} NumPy and PyTorch seeds are set to 42 for
	reproducibility.
\end{itemize}

\section{Inference Protocol}

Given a new split (test or holdout):
\begin{enumerate}
	\item Compute $\hat{y}^{\text{CB}}_{i,t}$ using the CatBoost model.
	\item Build sequences $\mathcal{S}_{i,t}$ when at least $L$ months of
	history exist; otherwise no CNN residual is produced for that row.
	\item Predict residuals $\hat{e}_{i,t}$ with the CNN and set
	$\hat{e}_{i,t}=0$ when no sequence is available.
	\item Return the final prediction
	$\hat{y}_{i,t} = \hat{y}^{\text{CB}}_{i,t} + \hat{e}_{i,t}$.
\end{enumerate}
This inference path performs no parameter updates.

\section{Hyperparameters Used}

\subsection*{CatBoost}
\begin{center}
\begin{tabular}{ll}
\toprule
Parameter & Value \\
\midrule
Loss / Eval metric & RMSE \\
Iterations & 1500 \\
Depth & 10 \\
Learning rate & 0.05 \\
L2 leaf reg & 3.0 \\
Random seed & 42 \\
Verbose & 100 \\
Task type & GPU (CPU fallback) \\
\bottomrule
\end{tabular}
\end{center}

\subsection*{Residual CNN}
\begin{center}
\begin{tabular}{ll}
\toprule
Parameter & Value \\
\midrule
Sequence length $L$ & 6 \\
Batch size & 1024 \\
Epochs & 30 \\
Learning rate & $10^{-3}$ \\
Weight decay & $10^{-4}$ \\
Optimizer & Adam \\
Training split & Validation only \\
\bottomrule
\end{tabular}
\end{center}

\appendix
\section{CatBoost from First Principles}

CatBoost is a gradient boosted decision tree method designed to handle
categorical features and reduce prediction bias. Let $\mathcal{L}(y, F(x))$ be a
loss function and define the additive model
\[
F_M(x) = \sum_{m=1}^{M} \eta_m h_m(x),
\]
where $h_m$ is a decision tree and $\eta_m$ is a learning rate. Gradient boosting
fits $h_m$ to the negative gradient of the loss:
\[
g_m(x_i) = -\left.\frac{\partial \mathcal{L}(y_i, F(x_i))}{\partial F}\right|_{F=F_{m-1}},
\quad h_m \approx \arg\min_h \sum_i (g_m(x_i) - h(x_i))^2.
\]
The model updates as $F_m(x) = F_{m-1}(x) + \eta_m h_m(x)$.

\subsection*{Oblivious (Symmetric) Trees}

CatBoost uses symmetric trees where the same splitting feature and threshold are
applied at each level. A depth-$d$ oblivious tree has $2^d$ leaves and can be
written as
\[
h(x) = w_{\ell(x)}, \quad \ell(x) \in \{0,\dots,2^d-1\},
\]
where $\ell(x)$ is determined by the sequence of binary tests along the tree.
Symmetry reduces overfitting and yields fast evaluation.

\subsection*{Categorical Features and Target Statistics}

For a categorical feature $c$, CatBoost builds target statistics (TS), such as
\[
\text{TS}(c=v) = \mathbb{E}[y \mid c=v],
\]
possibly combined with prior smoothing. To avoid target leakage, CatBoost uses
\emph{ordered boosting}: a random permutation of the dataset is fixed, and the TS
for each row is computed using only preceding rows in the permutation. This
makes the TS causal with respect to the ordered data and reduces bias.

\subsection*{Ordered Boosting}

Standard boosting can suffer from prediction shift when the same data are used
to compute target statistics and fit the model. CatBoost mitigates this by
building multiple permutations and averaging gradients across them. The result
is a lower-bias estimator while preserving the efficiency of tree boosting.

\subsection*{Regularization}

CatBoost uses $\ell_2$ regularization on leaf values and early stopping based on
the validation set. In this pipeline, the hyperparameter
\texttt{l2\_leaf\_reg} controls the strength of the leaf penalty, and the best
iteration is selected using the validation loss.

\end{document}
